\documentclass[11pt]{article}
\usepackage[T1]{fontenc}
\usepackage[utf8]{inputenc}
\usepackage{amsmath}
\usepackage{amssymb}
\usepackage{booktabs}
\usepackage[margin=1in]{geometry}
\usepackage[hidelinks]{hyperref}
\usepackage{graphicx}
\graphicspath{{figs/}}

\newcommand{\backend}[1]{\texttt{#1}}

\title{\bf Auditing Actual Execution and Output Portability in
Quantized LLM Inference:\\ A Measurement Study of llama.cpp}
\author{Anonymous}
\date{October 2026 (technical note)}

\begin{document}
\maketitle

\begin{abstract}
We measured whether the same quantized GGUF model, the same prompt, and
fully deterministic greedy decoding (temperature 0) produce the same
output across the compute backends of llama.cpp (release b11327,
official binaries, SHA256-recorded) on a 12-machine multi-vendor fleet
(NVIDIA, AMD, Intel; Linux and Windows). Auditing \emph{which device
actually executed each run} --- via per-prompt wall-time ratios against
the same-box CPU build --- turned out to be the decisive methodological
step: three distinct failure classes silently served inference from the
CPU while presenting as GPU backends, including the official ROCm build
on four AMD hosts (its backend library requires ROCm userspace SONAMEs
one major version newer than stock installs, and no Linux runtime is
bundled), a CUDA-13 build deprived of its runtime libraries by our own
launch environment, and a virtualized host with no GPU passthrough at
all. On runs verified to execute on the GPU, byte-exact agreement with
the same-box CPU reference is quantization-dominated and
monotone --- F16 88.0\%, Q8\_0 35.3\%, Q4\_K\_M 18.0\% of 8{,}670
prompt comparisons --- and nearly identical rates appear between
\emph{different builds' CPU implementations} on the same host (F16
92.7\%, Q8\_0 40.5\%, Q4\_K\_M 17.7\%), and even between the same CPU
build on different hosts. Divergence is concentrated at near-ties: at
all 7{,}875 first divergences, the alternative token lies within the
reference's recorded top-5 in 99.89\% of events (median logprob gap
0.062 nats, p90 0.231). Identical stacks are perfectly deterministic:
100\% agreement across two different physical machines with identical
GPU/driver/binary, and 100\% across a same-machine 12-hour re-run of
51 cells. We release the harness, per-token logprob records, and the
device-truth audit method. Practical conclusions: verify actual
execution placement before labeling any measurement with a backend
name; treat quantized artifacts as non-portable across stacks at
Q4-class precision; and do not use default-tool perplexity as a
conformance signal.
\end{abstract}

\section{Introduction}\label{sec:intro}

A greedy-decoded output is usually treated as a property of the model
and the prompt. This note reports a measurement study of the other
variable: the numerical stack that executes the model. We ask two
questions. First, \emph{does the intended backend actually execute the
workload}? Second, \emph{when the stack differs, how much does greedy
output differ, and where}?

The first question became the pivot of the study. We began by comparing
backends on a multi-vendor fleet and initially read the results as
backend behavior. A device-truth audit --- comparing every run's
per-prompt generation wall time against the same-box CPU build ---
revealed that three of our ``GPU backend'' columns had silently served
from the CPU: the official ROCm build on all four AMD hosts, the
CUDA-13 build on both hosts offering it, and CUDA on a virtualized
host. Only the first is an upstream packaging gap; the second was a
launch-environment mistake in our own harness (which had passed the
runtime-library path to device enumeration but not to the server
process); the third was a virtual machine with no GPU passthrough. All
three served HTTP requests normally, with no error on the console, and
two of the three even enumerated a device under \texttt{--list-devices}
--- enumeration, we conclude, is not evidence of execution.

The audit changed the study's claims in both directions. It killed
findings we had briefly believed (a ``CUDA-toolkit changes everything''
comparison; a ``perplexity defaults differ between tools'' hypothesis)
because those comparisons had in fact measured CPU implementations
against each other. It also sharpened the findings that survive: the
conformance structure of runs \emph{verified} to differ in execution,
the near-tie anatomy of divergence, and the perfect determinism of
identical stacks.

\section{Method}\label{sec:method}

\textbf{Fleet.} 12 physical machines (13 run labels; one machine was
measured twice, 12 hours apart, and its re-run is used only as a
stability control): NVIDIA (RTX 5060 Ti, RTX PRO 6000 Blackwell, TITAN
RTX, RTX 3060 Ti/Windows), AMD (RX 7900 XTX $\times$2, Radeon AI PRO
R9700, Strix Halo 8060S), Intel (Arc A770 Linux + Windows, Arc B60);
hosts span several CPU generations and both desktop OSes.

\textbf{Engine and models.} llama.cpp release b11327, official
prebuilt binaries only, SHA256 recorded. 17 core GGUF models:
Qwen2.5-Instruct 0.5B/1.5B/3B and SmolLM2-Instruct 135M/360M/1.7B at
Q4\_K\_M, Q8\_0, and F16/FP16 (official Qwen-org quantizations;
community re-quantizations for SmolLM2 --- provenance is a caveat for
family comparisons).

\textbf{Protocol.} 34 fixed raw-completion prompts (no chat template;
four categories), greedy decoding (temperature 0), 128-token cap,
prompt caching disabled, one request at a time. The server protocol
records, for every generated token on \emph{every} backend, the chosen
id, its logprob, and the top-5 ids with logprobs. Perplexity is
measured by the same release's \texttt{llama-perplexity} on a fixed
$\sim$5k-token corpus. The same-box CPU build is the reference for all
per-cell comparisons.

\textbf{Device-truth audit.} For every (machine, backend) cell we
compute the median over models of the ratio of median per-prompt
generation wall time to the same-box CPU build's. Observed ratios
bifurcate cleanly ($\approx 1.0$ vs $\le 0.67$); we verified a sample
of cells at the $\approx 1.0$ end with process maps, VRAM readings,
and library-load checks (\S\ref{sec:fallback}). Cells at
$\approx 1.0$ are classified \emph{CPU-fallback} and analyzed as a
separate population, not as GPU backends.

\section{Finding 1: silent CPU fallback}\label{sec:fallback}

Three failure classes, one user-visible behavior: the server accepts
requests and generates normally; nothing on the console distinguishes
it from a GPU deployment.

\textbf{(a) Missing backend userspace (upstream-relevant).} The
official ROCm build's \texttt{libggml-hip.so} requires
\texttt{libamdhip64.so.7}, \texttt{libhipblas.so.3}, and
\texttt{librocblas.so.5}. Stock Ubuntu hosts with ROCm 6.3.0 provide
\texttt{.so.6/.so.2/.so.4}; the release bundles no Linux ROCm runtime
(a Windows runtime zip exists; the Linux equivalent does not). The
backend library therefore never loads on any of our four AMD hosts ---
including an RX 7900 XTX, an officially supported target ---
\texttt{--list-devices} prints \texttt{(none)}, and the server runs
fully on the CPU, \emph{including with an explicit \texttt{-ngl 999}}:
process maps show no HIP library, VRAM stays idle, and generation
proceeds at exactly CPU-build speed. The perplexity tool is subject to
the same backend failure, so its values are host-CPU values; on our
data the ``ROCm'' perplexity matched the CPU build's in 68 of 68
cells, and every ``ROCm'' generation was bit-identical to the CPU
build's on the same host. Related reports exist upstream for Windows
and other setups; our contribution is the Linux SONAME detail, the
cross-host confirmation, and the benchmarking consequence: quality
measurements taken this way silently measure the CPU.

\textbf{(b) Runtime-library environment mistakes (our own).} Our
harness passed the CUDA runtime-library path
(\texttt{LD\_LIBRARY\_PATH}) to device enumeration and to the
perplexity tool, but not to the server process. On hosts where
\texttt{libcudart.so.13} exists only in the release's runtime tarball,
the CUDA-13 backend loaded for \texttt{--list-devices} --- which
therefore reported the GPU --- but not for the server, which silently
served from the CPU. An on-box A/B check isolates the effect: with the
environment set, the CUDA-13 build offloads under the default
\texttt{-ngl auto} at 667 tok/s; without it, even \texttt{-ngl 999}
runs on the CPU at 203 tok/s with 2\,MiB VRAM in use. The CUDA-12.8
build was unaffected only because those hosts carry system CUDA 12.9
libraries. We retract the preliminary ``CUDA-13 behaves like a
different backend'' observation that this artifact produced, and note
the general lesson: \texttt{--list-devices} output reflects the
environment of whatever process ran it, which is not necessarily the
environment of the server.

\textbf{(c) Absent hardware.} One host, labeled as a virtualized RTX
4090, has no NVIDIA device on its guest PCI bus at all (no
\texttt{/dev/nvidia*}, no driver library, only an Intel iGPU on
\texttt{lspci}). Every backend on it was CPU execution by
construction. We retain this host's cells only within the
CPU-fallback population and note that our earlier ``virtualized
4090'' hardware label was wrong.

\textbf{Detection.} All three classes were invisible to
\texttt{--list-devices} in at least one instance, and none produced a
startup error. The wall-time-ratio audit caught all three, and it is
cheap: it requires only a CPU-build baseline on the same host. We
recommend it as a standing check in any benchmarking or conformance
setup, alongside post-launch verification that the backend library is
actually mapped in the server process.

\section{Finding 2: conformance on verified runs}\label{sec:conformance}

On runs verified to execute on the GPU (CUDA on 4 machines, Vulkan on
10; 255 cells after removing the re-run), byte-exact agreement with
the same-box CPU reference is:

\begin{center}
\small
\begin{tabular}{lccc}
\toprule
Population & Q4\_K\_M & Q8\_0 & F16/FP16 \\
\midrule
GPU vs same-box CPU (n: 3060/3060/2550) & 18.0\% & 35.3\% & 88.0\% \\
Other builds' CPU impl.\ vs same-box CPU (n: 1428/1428/1190) & 17.7\% & 40.5\% & 92.7\% \\
Same CPU build, different host (per pair) & 19.6--20.1\% & 39.2--40.7\% & 92.9--94.7\% \\
\bottomrule
\end{tabular}
\end{center}

Three observations. First, quantization dominates: within every
population the rates are monotone in precision, and the pooled
``$\sim$45--48\% agreement'' that a fixed quant mix produces is a
mixture artifact --- per-precision rates span 18\% to 93\%. Second,
the comparison type is a second-order effect: GPU-vs-CPU, one
build's CPU implementation vs another's (including the ROCm and
CUDA-13 fallback columns, reclassified), and even the same CPU build
on different host CPUs all land within a few points of each other at
each precision. Byte-exact portability across \emph{any} numerical
change of stack is essentially gone at Q4-class precision and
partially preserved at F16. Third, the CPU reference is itself
stack-relative: different hosts running the same CPU build agree with
each other no more than GPU backends agree with their local CPU, and
CPU outputs form bit-identical equivalence classes across subsets of
our hosts. ``Same result as CPU'' therefore means ``same result as
this host's CPU build''.

\section{Finding 3: identical stacks are deterministic}\label{sec:determinism}

Two controls bound the divergence findings from above. Two physically
distinct machines with the same GPU (RX 7900 XTX), same driver, same
OS, and same Vulkan binary produced bit-identical generations for
578/578 prompt comparisons --- cross-machine determinism is perfect
when the stack is identical. The same machine re-run 12 hours later
(51 cells, 1{,}734 prompts; CUDA, CUDA-13, and Vulkan) reproduced its
earlier outputs bit-for-bit in 100.0\% of comparisons. Divergence in
\S\ref{sec:conformance} is therefore a property of differing stacks,
not of nondeterminism within a stack. (We separately note that the
engine's own documentation warns that prompt-cache reuse can
introduce run-to-run nondeterminism via batch-shape-dependent
prefill/decode arithmetic; our protocol runs with caching disabled.)

\section{Finding 4: divergence anatomy}\label{sec:mechanism}

At every first-divergence position (7{,}875 events; 6{,}903 after
removing identical-stack duplicates), the diverging backend's chosen
token lies within the reference's recorded top-5 in 99.89\% of events
(9 exceptions, all in one SmolLM2 Q4\_K\_M cell family). Because both
sides' top-5 logprobs are recorded, the gap can be measured in both
distributions: the reference's preference for its own argmax over the
backend's choice has median 0.062 nats (p90 0.231, p99 0.564); the
backend's preference for its choice over the reference's has median
0.059 nats (p90 0.226). Hand inspection of sampled divergences shows
rank-1$\leftrightarrow$2 swaps. Divergence at greedy decoding is
therefore concentrated at near-ties --- compatible with the
rounding-order mechanism assumed in prior work, though our data does
not by itself exclude small systematic implementation differences
that would also preferentially flip low-margin decisions.

\section{Related work}\label{sec:related}

Cross-platform numerical deviation in DNN inference, with platform
equivalence classes, was established for vision models by Schl\"ogl
et al.\ (NeurIPS 2023). He (Thinking Machines Lab, 2025) identifies
batch-shape dependence as the root cause of temperature-0
nondeterminism and ships batch-invariant kernels; the engine's own
cache-reuse warning is an instance of this class. Yuan et al.\ (2025,
v2: ``Understanding and Mitigating Numerical Sources of
Nondeterminism in LLM Inference'') measure greedy divergence across
GPU type, count, batch size and precision at bf16 scale and propose
LayerCast; Du et al.\ (TMLR 2026) measure cross-precision greedy
divergence with margin-thresholded mitigation; Dutta et al.\
(NeurIPS 2024) show aggregate accuracy metrics hide flips in
compressed models --- our perplexity observations are an instance.
Pape et al.\ (2026) quantify backend-choice effects across engines
including llama.cpp. Deterministic or verified inference is addressed
by Gond et al.\ (LLM-42, 2026), Cooper et al.\ (2026, fixed-reduction
GEMMs), and TOPLOC (Ong et al., 2025) for verifiable inference. In
software engineering, CRADLE (ICSE 2019) cross-validated deep-learning
library backends; our harness follows that tradition for
autoregressive LLM inference, and its device-truth audit extends it
with an execution-placement check. Training-side tooling variance is
covered by Zhuang et al.\ (2021) and Pham et al.\ (ASE 2020).

\section{Limitations}\label{sec:limitations}

One engine, one release (b11327) --- rates will drift with any
release, while the silent-fallback classes and the audit method are
the durable content. Models are small (135M--3B) by design (exact
byte-level comparison needs exactly reproducible runs; larger models
are future work). Raw-completion prompts without chat templates
raise continuation entropy relative to deployed chat use. The
near-tie evidence is compatible with, but does not uniquely
establish, rounding-order causation. Perplexity comparisons across
backends are confounded by the fallback classes and by
prefill-vs-decode kernel differences. SYCL contributed no complete
cells (Linux binaries failed to load the oneAPI runtime; the Windows
build loaded, demonstrably offloaded, then crashed mid-run).

\section{Conclusion}\label{sec:conclusion}

For practitioners: (1) verify what executed --- the wall-time ratio
against a CPU baseline is a cheap, reliable detector of silent CPU
fallback, and \texttt{--list-devices} is not; (2) expect byte-exact
greedy output to be non-portable across any stack change at
Q4-class precision (18\% agreement), partially portable at Q8
(35\%), and mostly preserved at F16 (88\%); (3) identical stacks are
perfectly deterministic --- pin backend, binary, GPU, and driver if
you need reproducibility; (4) do not treat default-tool perplexity
as evidence of behavioral equivalence. We release the harness, the
per-token logprob dataset, and the audit tooling.

\end{document}
